\documentclass[10pt,a4paper]{amsart}
\usepackage[margin=19mm]{geometry}
\usepackage{amsmath,amssymb,mathtools}
\usepackage[scr=rsfs,cal=euler]{mathalfa}
\usepackage{xfrac}
\usepackage[unicode,hidelinks,bookmarks=false]{hyperref}
\newcommand{\ie}{i.e.\ }
\newcommand{\cf}{cf.\ }
\newcommand{\resp}{respectively\ }
\newcommand{\Subsection}{\subsection}
\providecommand{\nobreakdash}{\nobreak\hbox{-}}
\setlength{\emergencystretch}{2em}
\multlinegap0pt
\usepackage{amssymb}
\usepackage{mathtools}
\usepackage[shortlabels]{enumitem}
\newtheorem{theorem}{Theorem}
\newtheorem{lemma}{Lemma}
\usepackage{cleveref}
\crefname{theorem}{Theorem}{Theorems}
\crefname{lemma}{Lemma}{Lemmas}
\newcommand{\E}{\mathbb E}
\newcommand{\one}{\mathbf 1}
\newcommand{\splus}{s^{+}}
\title{A Vertex-Localized Positive Square-Energy Strengthening\\ of Tur\'an's Theorem}
\author{Abhay Jayarajan, M. Rajesh Kannan, Shivaramakrishna Pragada, Rahul Roy}
\date{Source: arXiv:2608.24861v2 (CC BY 4.0)}
\begin{document}
\maketitle

\noindent\emph{Source and licence.} A Vertex-Localized Positive Square-Energy Strengthening\\ of Tur\'an's Theorem. Version arXiv:2608.24861v2 (CC BY 4.0), distributed by arXiv under the Creative Commons Attribution 4.0 International licence, http://creativecommons.org/licenses/by/4.0/, as recorded in the arXiv licence metadata for this exact version and shown on the abstract page https://arxiv.org/abs/2608.24861v2. Full licence notice: \texttt{licenses/arxiv-2608.24861-CC-BY-4.0.txt}.\par\smallskip

\noindent\textit{Editorial scope.} Counted unit: the article's lemma lem:local-clique-rigidity (source0.tex lines 695--704). The article's complete original proof of it is delivered with it (source0.tex lines 706--819, one \texttt{proof} environment of 114 lines). No mathematics is written, repaired or re-derived: the statement and the proof are the article's own lines, in the article's own notation, and no retained line is substituted or re-flowed.  Retained uncounted as the source-local context the proof needs: lines 404--417; lines 446--453; lines 456--469; lines 475--481; lines 671--673; lines 688--690.  Nothing was omitted from any retained span, so the package carries no omission record.  No retained line contains a citation, so the wrapper carries no bibliography and no bibliography entry.  No retained line contains a figure environment, an image-inclusion command or drawing code, so no image is delivered and the package declares no asset dependency.  Editorial formatting only, in addition: an AMS \texttt{amsart} preamble that restates the article's own declarations and macros used by the retained lines, namely the environments \texttt{theorem}, \texttt{lemma} on separate counters, together with the standard packages those lines need.  The retained environments print in this document as Theorem 1, Lemma 1 in order of appearance, whereas the article prints the same items with the numbering of its own full document.  The complete native source of the article as distributed in the arXiv e-print for this exact version is bundled unchanged as \texttt{sources/208523/source0.tex}.
	\begin{theorem}
		\label{thm:localized-DNN}
		Let $G$ be a finite simple graph and let $M$ be a doubly nonnegative
		matrix indexed by $V$.  Set
		\[
		T:=\one^{\top}M\one,
		\qquad
		S:=\sum_{uv\in E}\sqrt{M_{uv}}.
		\]
		Then
		\begin{equation}\label{eq:localized-DNN}
			2S\leq \sum_{v\in V}\left(1-\frac1{c(v)}\right)\sqrt{T}.
		\end{equation}
	\end{theorem}

		\begin{equation}\label{eq:Theta-bound}
			\Theta
			:=
			\sum_{v\in V}\tau_v^2p_v
			=
			\E\left(\sum_{v\in K}\tau_v^2\right)
			\leq\eta.
		\end{equation}

		\begin{align}
			D^2
			&=
			\left(
			\sum_{v\in V}
			\tau_v\sqrt{p_v}\,\frac1{\sqrt{p_v}}
			\right)^2 \notag\\
			&\leq
			\left(\sum_{v\in V}\tau_v^2p_v\right)
			\left(\sum_{v\in V}\frac1{p_v}\right)
			=\Theta B
			\leq\eta B.
			\label{eq:D-square}
		\end{align}

		\begin{equation}\label{eq:H-upper}
			H
			\leq
			\frac12\left(n^2-B\right)
			\leq
			\frac12\left(n^2-\frac{D^2}{\eta}\right).
		\end{equation}

	\begin{equation}\label{eq:connected-equality}
		\sqrt{\splus(G)}=\sum_{v\in V}\left(1-\frac1{c(v)}\right).
	\end{equation}

	\begin{equation}\label{eq:DNN-equality-special}
		2S=\sum_{v\in V}\left(1-\frac1{c(v)}\right)\sqrt{T}.
	\end{equation}

	\begin{lemma}
		
		\label{lem:local-clique-rigidity}
		Under~\eqref{eq:connected-equality},
		\[
		c(v)=\omega(G)
		\qquad\text{for every }v\in V.
		\]
		Moreover, every Caro-Wei ordering has exactly $\omega(G)$ pivots.
	\end{lemma}

	\begin{proof}
		Retain all notation from the proof of \Cref{thm:localized-DNN}:
		\[
		\tau_v=\frac1{c(v)},
		\quad
		D=\sum_v\tau_v,
		\quad
		\eta=\E\left(\frac1R\right),
		\]
		\[
		B=\sum_v\frac1{p_v},
		\quad
		H=\sum_{uv\in E}\frac1{q_{uv}},
		\quad
		\Theta=\sum_v\tau_v^2p_v,
		\quad
		\mathcal W=\sum_{uv\in E}M_{uv}q_{uv}.
		\]
		Define
		\begin{equation}\label{eq:upper-H-W}
			\widehat H
			:=\frac12\left(n^2-\frac{D^2}{\eta}\right),
			\qquad
			\widehat{\mathcal W}
			:=\frac12(1-\eta)T.
		\end{equation}
		The proof of \Cref{thm:localized-DNN} gives
		\begin{equation}\label{eq:equality-long-chain}
			S^2
			\leq\mathcal W H
			\leq\widehat{\mathcal W}\,\widehat H
			\leq\frac{(n-D)^2T}{4}.
		\end{equation}
		By~\eqref{eq:DNN-equality-special}, the first and last terms in
		\eqref{eq:equality-long-chain} are equal.
		
		Because $G$ is connected and nonempty, every first pivot has a neighbor,
		so every outcome has $R\geq2$.  Hence $\eta<1$ and
		$\widehat{\mathcal W}>0$.  Also $H>0$, since $E\neq\varnothing$.
		Since
		\[
		0\leq\mathcal W\leq\widehat{\mathcal W},
		\qquad
		0<H\leq\widehat H,
		\]
		equality of the endpoint products in~\eqref{eq:equality-long-chain}
		forces
		\begin{equation}\label{eq:W-H-equality}
			H=\widehat H.
		\end{equation}
		
		The proof of~\eqref{eq:H-upper} gives the more detailed chain
		\[
		H
		\leq\frac12(n^2-B)
		\leq\frac12\left(n^2-\frac{D^2}{\eta}\right)
		=\widehat H.
		\]
		Together with~\eqref{eq:W-H-equality}, this yields
		\begin{equation}\label{eq:B-exact}
			B=\frac{D^2}{\eta}.
		\end{equation}
		Now~\eqref{eq:D-square}, \eqref{eq:Theta-bound}, and
		\eqref{eq:B-exact} imply
		\[
		D^2\leq\Theta B\leq\eta B=D^2.
		\]
		Since $B>0$, equality holds throughout and therefore
		\begin{equation}\label{eq:Theta-eta}
			\Theta=\eta.
		\end{equation}
		
		For a fixed total ordering $\pi$, let $K_\pi$ be its pivot set and
		$R_\pi:=|K_\pi|$.  Pointwise,
		\begin{equation}\label{eq:pointwise-tau}
			\sum_{v\in K_\pi}\tau_v^2\leq\frac1{R_\pi}.
		\end{equation}
		The expectation of the difference between the two sides is
		$\eta-\Theta=0$ by~\eqref{eq:Theta-eta}.  Every ordering has positive
		probability, and the difference in~\eqref{eq:pointwise-tau} is
		nonnegative.  Hence equality holds in~\eqref{eq:pointwise-tau} for every
		ordering.
		
		For each pivot $v\in K_\pi$, we have $\tau_v\leq1/R_\pi$.  There are
		$R_\pi$ pivots, and
		\[
		\sum_{v\in K_\pi}\tau_v^2
		=\frac1{R_\pi}
		=R_\pi\left(\frac1{R_\pi}\right)^2.
		\]
		% Thus every summand attains its individual maximum:
		Thus we must have
		\begin{equation}\label{eq:c-pivot-R}
			c(v)=R_\pi
			\qquad
			(v\in K_\pi)
		\end{equation}
		for every ordering $\pi$.
		
		Let $uv\in E$.  Choose an ordering whose first two vertices are $u$
		and $v$, in that order.  Then $u$ is the first pivot and $v$ is the
		second pivot, because $v\in N(u)$ and is first in the induced order on
		$N(u)$.  Applying~\eqref{eq:c-pivot-R} to both pivots gives
		$c(u)=c(v)$.  Since $G$ is connected, $c(v)$ is constant over all
		vertices.  A vertex belonging to a maximum clique has local clique number
		$\omega(G)$, so
		\[
		c(v)=\omega(G)
		\qquad(v\in V).
		\]
		Finally, in every ordering the first pivot satisfies
		$c(x_0)=R_\pi$ by~\eqref{eq:c-pivot-R}; hence
		$R_\pi=\omega(G)$ for every ordering.
	\end{proof}

\end{document}
